% concatenated from mainV2.tex

\documentclass[a4paper, 12pt]{amsart}
\usepackage{amssymb, amsfonts, amsmath, bm}
\usepackage[all]{xy}
\usepackage[shortlabels]{enumitem}
\usepackage{hyperref, color}
\usepackage{commath}
\usepackage{esdiff}
\usepackage{graphicx,subfigure}
\usepackage[center]{caption}
\usepackage{amsthm}
\usepackage{diagbox}
\usepackage{multirow}
\usepackage{esdiff}
\usepackage[pagewise]{lineno}
\usepackage[utf8]{inputenc}
\usepackage{soul}
\usepackage{comment} 
\usepackage{pifont}
\usepackage{hyperref, color, xcolor}
\hypersetup{
	colorlinks = true,
	linkcolor = blue,
	filecolor = blue,
	urlcolor = red,
	citecolor = red
}
\usepackage{cite}
\usepackage{lineno}
%\linenumbers % Ativa a numeração de linhas
%\renewcommand{\qedsymbol}{\textcolor{blue}{$\blacksquare$}}
%\renewcommand{\qedsymbol}{\textcolor{orange}{$\lozenge$}}  % losango
%\renewcommand{\qedsymbol}{\textcolor{orange}{$\blacksquare$}}
%\renewcommand{\qedsymbol}{\textcolor{orange}{$\blacklozenge$}}
\renewcommand{\qedsymbol}{\textcolor{orange}{\ding{113}}}  % estrela sólida roxa
\newcommand{\inner}[2]{\left\langle #1,#2\right\rangle}
%\renewcommand{\qedsymbol}{\textcolor{purple}{$\diamondsuit$}}  % naipe


\usepackage[nameinlink,capitalise]{cleveref} % References show what kind of thing is being referenced (theorem, lemma, proposition...)

\renewcommand{\fullwidthdisplay}{}

\usepackage[a4paper,twoside,top=1in, bottom=1in, left=0.8in, right=0.8in]{geometry}

%\makeatletter
%\renewcommand{\author@andify}{
	%	\nxandlist {\unskip ,\penalty-1 \space\ignorespaces}
	%	{\unskip {} \@@and~}
	%	{\unskip \penalty-2 \space \@@and~}
	%}
%\makeatother
\DeclareMathOperator{\Ric}{Ric}
\DeclareMathOperator{\divv}{div}

\newcommand{\Rbb}{\mathbb{R}}
\newcommand{\Sbb}{\mathbb{S}}

\newcommand{\tg}{\tilde{g}}
\newcommand{\tX}{\widetilde{X}}
\newcommand{\talpha}{\tilde{\alpha}}
\newcommand{\tv}{\tilde{v}}
\newcommand{\tV}{\widetilde{V}}
\newcommand{\tnabla}{\widetilde{\nabla}}
\newcommand{\tH}{\widetilde{H}}
\newcommand{\tdelta}{\tilde{\delta}}
\newcommand{\ta}{\tilde{a}}
\newcommand{\tmu}{\tilde{\mu}}
\newcommand{\tE}{\tilde{E}}

\newcommand{\bOmega}{\bar{\Omega}}

\newtheorem{theorem}{Theorem}[]
\newtheorem{lemma}[theorem]{Lemma}
\newtheorem{proposition}[theorem]{Proposition}
\newtheorem{corollary}[theorem]{Corollary}
\newtheorem*{conjecture*}{Conjecture}
%\newtheorem{question}[theorem]{Question}

\newtheorem{theo}{Theorem}
\renewcommand{\thetheo}{\Alph{theo}}
\newtheorem{question}[theo]{Question}
\newtheorem{rem}[theo]{Remark}
\newtheorem{cor}[theo]{Corollary}
\newtheorem{prop}[theo]{Proposition}


\newtheorem{claim}{Claim}

\theoremstyle{definition}
\newtheorem{example}{Example}
\newtheorem*{remark*}{Remark}

\newtheorem{remark}[theorem]{Remark}
\newtheorem*{acknowledgements}{Acknowledgements}
\newtheorem{definition}[theorem]{Definition}
\newtheorem{item1}{(1)}

\numberwithin{figure}{section}

\usepackage{graphicx} 
\title[On Generalized $m$-Quasi-Einstein Manifolds]{On Generalized $m$-Quasi-Einstein Manifolds}

\author[Alcides de Carvalho, Anderson Lima and W. O. Costa-Filho]{Alcides de Carvalho, Anderson Lima and W. O. Costa-Filho}

        \address{Universidade Federal de Pernambuco\\
		Departamento de Matematica\\
		 Av. Jorn. Aníbal Fernandes, s/n, Cidade Universitária, Recife - PE, Brazil, CEP:  50740-560.}	
	\email{alcides.junior@ufpe.br}
  \address{
  CPMAT - IM, Universidade Federal de Alagoas, Maceió, AL, 57072-970, Brazil}
\email{jose.lima@im.ufal.br}
  
\address{
  Universidade Federal de Alagoas, Campus Arapiraca, Av. Manoel Severino Barbosa, S/N, Bom Sucesso, 
 Arapiraca - AL, Brazil, CEP: 57309-005.}
\email{wagner.filho@arapiraca.ufal.br}
%0000-0003-0034-2580
%0000-0003-2568-6590
%0000-0002-9536-5322

\date{\today}
 
\begin{document}
\keywords{m-Quasi-Einstein manifolds · Killing vector fields · Scalar curvature}

\subjclass{Primary 53C25; Secondary 53C21, 53C65}
\maketitle

\begin{abstract}
In this paper, we study generalized $m$-quasi-Einstein manifolds $(M^n,g,X,\lambda)$ under suitable constraints on the potential vector field $X$. We establish that, under suitable integral assumptions, the potential vector field is a Killing vector field, extending earlier results of Sharma \cite{Rsharma2025} to the generalized setting. In addition, we show that if $X$ is divergence-free, then $X$ is a Killing vector field. We derive consequences under sign conditions on $m$ and $\lambda$, including triviality results. We also revisit a recent theorem by Ghosh \cite{ghosh}, discuss a subtle issue in the argument, and provide a new formulation and proof. Finally, we obtain rigidity results when $X$ is a geodesic vector field.
\end{abstract}


\section{Introduction}

Besides their geometrical importance, Riemannian manifolds admitting an Einstein-like structure are of interest in Mathematical Physics where they are usually referred as quasi-Einstein metrics. In recent years, several generalizations of this notion have been introduced and extensively investigated in the literature. Among them, the class of generalized m-quasi-Einstein manifolds has attracted considerable attention due to its rich geometric nature and its connections with Ricci solitons, warped product Einstein metrics and many gravitational contexts in spacetime. See, for instance, Bahuaud, Gunasekaran, Kunduri, and Woolgar in \cite{Bahuaud2024}, as well as Wylie in \cite{Wylie2023} and the relevant references therein.

An $n$-dimensional Riemannian manifold $(M^n, g)$, with $n \geq 2$, is said to be a \emph{generalized $m$-quasi-Einstein manifold} for a nonzero constant $m \in \mathbb{R}$ if there exists a smooth vector field $X$ on $M^n$ and a smooth function $\lambda \colon M \to \mathbb{R}$ such that the following equation holds:
\begin{equation}\label{eq:quasi-einstein}
\operatorname{Ric} + \frac{1}{2} \mathcal{L}_X g - \frac{1}{m} X^\flat \otimes X^\flat = \lambda g,
\end{equation}
where $\operatorname{Ric}$ denotes the Ricci tensor, $\mathcal{L}_X g$ is the Lie derivative of the metric $g$ in the direction of $X$, $X^\flat$ is the 1-form metrically dual to $X$, and $\otimes$ denotes the tensor product. Throughout this paper, we adopt the convention that $(M^n,g,X,\lambda)$ is called an \emph{$m$-quasi-Einstein manifold} for solutions of \eqref{eq:quasi-einstein} with $\lambda$ constant.

This notion was introduced by Barros and Gomes in \cite{BarrosGomes} and was also considered by Barros and Ribeiro in \cite{BarrosRibeiro2014}, when $X$ is the gradient of a smooth potential function on $M^n.$ More recently, these equations have been referred to as the \emph{generalized extremal horizon equations}. See Definition 1.4 in Colling and Dunajski \cite{colling2025quasi}, where several related results are established. From the perspective of black hole theory, they arise as a generalization of the near-horizon geometry equations, which characterize the geometry of degenerate Killing horizons in spacetime (see also Lewandowski and Kaminski \cite{kaminski2024extreme}). 

Generalized $m$-quasi-Einstein manifolds are special generalizations of Ricci solitons, and a main motivation for studying this class of manifolds is their connection with general relativity. In addition, if $m$ is a positive integer, quasi-Einstein manifolds correspond as a base space to Einstein warped products. For more details on this subject, including some appropriate examples and important characterization results, we recommend that the reader consult \cite{BarrosGomes, BarrosRibeiro2014} and Ghosh in \cite{ghosh,Gho20}. An generalized $m$-quasi-Einstein manifold is called trivial when $X$ is identically zero. In what follows, $S$ denotes the scalar curvature of $M^n$, defined as the trace of $\operatorname{Ric}$.

Given their significance, a number of results have been derived with regard to quasi-Einstein manifolds under various assumptions. In \cite{bahuaud2022static} Bahuaud et al. studied the case where the dual one-form to $X$ is closed. In \cite{Bahuaud2024}, the authors considerably addressed the solutions of \eqref{eq:quasi-einstein} with the vector field $X$ divergence-free but not identically zero. In particular, they proved that if $M^n$ is closed and $m \neq -2$, then $X$ is a Killing vector field. As a consequence, the scalar curvature $S$ is constant.

More recently, Cochran \cite{Cochran2025} and Costa-Filho \cite{CostaFilho2024} independently showed that, on a closed $m$-quasi-Einstein manifold $(M^n,g,X,\lambda)$, the scalar curvature $S$ is constant if and only if $X$ is a Killing vector field. It is worth emphasizing that Cochran’s result requires the additional assumption $m \neq -2$. 
As pointed out by Cochran in \cite{Cochran2026}, this result was first shown by Ghosh in \cite{Gho20}. The implication that constant scalar curvature implies that $X$ is a Killing vector field is Theorem~4.2 of \cite{Gho20}. Furthermore, equation~(5.18) in \cite{Gho20} yields the converse implication.

Motivated by these developments, R. Sharma \cite{Rsharma2025} obtained a closely related rigidity result under a weaker hypothesis. More precisely, Theorem~1.1 in \cite{Rsharma2025} states that if $(M^n,g,X,\lambda)$ is a closed $m$-quasi-Einstein manifold, then $X$ being a Killing vector field implies that the scalar curvature $S$ is constant. Conversely, if
$ 
\int_M (\mathcal{L}_X S)\, dM\le 0,
$ in particular, if $S$ is constant along the flow of $X$ then $X$ must be a Killing vector field. At the time of publication, neither \cite{Cochran2025}, \cite{CostaFilho2024}, nor Sharma’s work were aware of Ghosh’s earlier result.


Inspired by Sharma’s result, we prove the following rigidity statement under analogous integral assumptions in the generalized $m$-quasi-Einstein setting.

\begin{theo}
Let $(M^n,g,X,\lambda)$ be a closed generalized $m$-quasi-Einstein manifold. Assume that
$$
\int_M \langle X,\nabla S\rangle\, dM \le 0
\quad \text{and} \quad
\int_M \langle X,\nabla \lambda\rangle\, dM \ge 0.
$$
Then $X$ is a Killing vector field.
\end{theo}
%\begin{theorem}
%Let $ (M^n, g, X,\lambda)$ be a closed generalized $ m $-quasi-Einstein manifold such that $\int_M \langle X,\nabla S\rangle\, dM \le 0.$ If $$\int_M \langle X,\nabla \lambda\rangle\, dM \ge 0,$$  then $X$ is a Killing vector field.
%\end{theorem}

\begin{remark}
If $\lambda$ is constant, then the second integral vanishes and the result reduces to the corresponding theorem of Sharma.
\end{remark}

Before continuing, we recall the following interesting and useful fact, which extends the previously mentioned results.

\begin{prop}\label{prop-div}
    Let $ (M^n, g, X,\lambda)$ be a closed generalized $ m $-quasi-Einstein manifold. If $X$ is divergence-free, then $X$ is a Killing vector field.
\end{prop}

In this scenario, very recently Ghosh~\cite{ghosh} established the following result.

\begin{theorem}[Ghosh, Theorem 6.1]
Let $(M^n, g, X, \lambda)$ be a closed generalized $m$-quasi-Einstein manifold. If $X$ is divergence-free, then $X$ is a Killing vector field of constant length and the scalar curvature $S$ is constant.
\end{theorem}
However, this statement is not valid in the generalized setting. Indeed, the proof in \cite{ghosh} relies on identity (5.3), which is derived under the assumption that $\lambda$ is constant. In the generalized framework, however, $\lambda$ is not assumed to be constant, and therefore the argument does not directly apply. A detailed discussion of this point is presented in Section~\ref{sec:proofs}.
Accordingly, the correct statement in the generalized setting is the following.

\begin{theo}\label{teo6.1}
Let $(M^n,g,X,\lambda)$ be a closed generalized $m$-quasi-Einstein manifold. If $X$ is divergence-free, then $X$ is a Killing vector field of constant length. Moreover, the scalar curvature $S$ is constant if and only if $\lambda$ is constant.% Furthermore, in dimension \(n=2\), the scalar curvature (S) is necessarily constant.
\end{theo}

The following example shows that the original statement of Theorem 6.1 in~\cite{ghosh} is false in the generalized setting.
\begin{example}
Consider the Riemannian manifold
$
(M^3,g)=\bigl((0,\infty)\times\mathbb{R}^2,\,dx^2+x^{-2}dy^2+x^{-4}dz^2\bigr),
$
the vector field
$
X=\partial_y,
$ on \(M\),
the real constant
$
m=\frac14,
$
and the smooth function
$
\lambda=-\frac{8}{x^2}$ of \(M\).
Since the metric coefficients do not depend on the variable \(y\), the vector field \(X\) is Killing. Moreover, a direct computation yields
$\Ric=-\frac{8}{x^2}\,dx^2-\frac{4}{x^4}\,dy^2-\frac{8}{x^6}\,dz^2,
$ and
$X^\flat\otimes X^\flat=x^{-4}\,dy^2.
$
Since \(\mathcal{L}_Xg=0\), we obtain
$
\Ric+\frac12\mathcal{L}_Xg-\frac1mX^\flat\otimes X^\flat
=-\frac{8}{x^2}\,g.
$
Therefore, \((M^3,g,X,\lambda)\) satisfies the generalized \(m\)-quasi-Einstein equation \eqref{eq:quasi-einstein}.
Furthermore, the scalar curvature is
$
S=-\frac{20}{x^2},
$ not constant.

%This example shows that generalized \(m\)-quasi-Einstein manifolds with Killing potential vector fields may have nonconstant scalar curvature.
\end{example}

A smooth vector field $Y$ on a Riemannian manifold $(M^n,g)$ is called a \textit{geodesic vector field} if $\nabla_YY=0$, where $\nabla$ is the Levi-Civita connection on $M^n$. Examples of geodesic vector fields include Killing vector fields of constant length on a Riemannian manifold. Recall that a vector field on $M^n$ is a geodesic vector field if each integral curve of $Y$ is a geodesic in $M^n.$ Physically, geodesic vector fields naturally appear in many important situations (see, for example, \cite{deshmukh2019geodesic,yano1961geodesic}).

Now, following Stepanov and Shandra in \cite{StepanovShandra2003}, we say that a smooth vector field $Y$ on $(M^n,g)$ is an \textit{infinitesimal harmonic transformation} if the one-parameter group of local transformations of $M^n$ generated by $Y$ consists of local harmonic diffeomorphisms.

As an immediate consequence of Lemma \ref{eq10}, we obtain the following improved version of Theorem 5.1 in \cite{ghosh}, to the generalized setting

\begin{theo}\label{D}
Let $(M^n,g,X,\lambda)$ be a generalized $m$-quasi-Einstein manifold with geodesic potential vector field $X$. Assume that $X$ is an infinitesimal harmonic transformation. If $n=2$, then $X$ is Killing. If $n\geq 3$, then $X$ is Killing if and only if $\lambda$ is invariant along the flow of $X$, that is,
$ 
X(\lambda)=0.
$
\end{theo}

%By Proposition~B and Theorem~C, the divergence-free condition implies that $X$ is a Killing vector field. In particular, $\lambda$ is constant, and therefore the results obtained under the divergence-free assumption naturally reduce to the classical $m$-quasi-Einstein setting.

In \cite{CostaFilho2024}, the third author presented a proof of Proposition 2.4 in \cite{Bahuaud2024} (see page 8) as a consequence of Lemma 1 in \cite{BarrosRibeiro2014}. In this work, we extend this approach to the generalized setting by establishing an analog of Lemma 1 from \cite{BarrosRibeiro2014} in this context (see Lemma~\ref{lem:laplacianX}). As an application, we derive the following fact.

\begin{cor}
Let $(M^n, g, X, \lambda)$ be a closed generalized $m$-quasi-Einstein manifold.
Assume that $X$ is divergence-free.
If $m < 0$ and $\lambda \leq 0$ on $M$, then $X$ vanishes identically.
\end{cor}




In the same work \cite{ghosh}, Ghosh proved that if $(M^n,g,X,\lambda)$ is a closed $m$-quasi-Einstein manifold with geodesic potential $X$ and non-positive Ricci tensor, then $X=0$ and $M^n$ is Ricci-flat, or $X=0$ and $M^n$ is Einstein, or
$
\lambda = -\frac{1}{m}|X|^2.
$
We extend this result from the $m$-quasi-Einstein setting to the class of generalized $m$-quasi-Einstein manifolds.

\begin{theo}
Let $(M^{n}, g, X, \lambda)$ be a closed generalized $m$-quasi-Einstein
manifold with geodesic potential vector field $X$. Assume that  $\int_M \operatorname{Ric}(X,X) dM \le 0$ and
$\int_M \langle X,\nabla \lambda\rangle\, dM \le 0$.
Then one of the following holds
\begin{itemize}
    \item[(i)] If \(\lambda =0\), then \(X=0\) and \((M^n,g)\) is Ricci-flat;

    \item[(ii)] If \(\lambda m>0\) on \(M^n\), then \(X=0\) and \((M^n,g)\) is Einstein;

    \item[(iii)] If \(X\not=0\), then
    $    \lambda=-\frac1m|X|^2
    $
    on the set \(\{X\neq0\}\).
\end{itemize}
\end{theo}

\begin{remark}
Under the assumptions of Proposition~\ref{prop:parallel} ($m  \geq 2(4-n)$), the vector field \(X\) is parallel and therefore \(|X|\) is constant. Hence, if \(X\not=0\), the relation
$
|X|^2=-m\lambda
$
implies that \(\lambda\) is constant and has sign opposite to that of \(m\).
\end{remark}
\section{Preliminaries}

In this section, we shall present some preliminaries which will be important for the establishment of the desired results. Throughout this paper, \((M^{n},g)\) denotes a connected and orientable Riemannian manifold.% Let $(M^n, g, X, \lambda)$ be a generalized $m$-quasi-Einstein manifold. 

We notice that taking the trace of equation \eqref{eq:quasi-einstein}, we arrive at  
\begin{equation}\label{eq2}
S + \operatorname{div}X = \frac{1}{m} |X|^2 + n\lambda.
\end{equation}


In what follows, we recall a key identity due to Barros and Ribeiro.
In the case of $m$-quasi-Einstein manifolds, that is, $\lambda$ is constant, it can be found in \cite[p.~215]{BarrosRibeiro2012}. Its extension to the generalized setting, with $X=\nabla f$, was obtained in \cite{BarrosRibeiro2014}. For the general case, this result was also proved by Ghosh in \cite[formula 6.5]{ghosh}. Here, $\Delta$ denotes the Laplacian operator on $M^n$.

\begin{lemma}\label{lem:laplacianX}
Let $(M^n,g,X,\lambda)$ be a generalized $m$-quasi-Einstein manifold. Then the following
identity holds
$$
\frac{1}{2}\,\Delta |X|^{2}
= |\nabla X|^{2}
- \mathrm{Ric}(X,X)
+ \frac{2}{m}\,|X|^{2}\,\mathrm{div}\,X
+ (2-n)\,\langle X,\nabla \lambda\rangle.
$$
\end{lemma}


When $\lambda$ is constant, the last term vanishes and the original formula is recovered. Lemma \ref{lem:laplacianX} yields the following result of rigidity under natural integral assumptions, which extends Theorem~1 of Barros and Ribeiro \cite{BarrosRibeiro2014} to the generalized setting.
\begin{corollary}
Let $(M^n,g,X,\lambda)$ be a closed generalized $m$-quasi-Einstein manifold with \\ $\int_M \langle X,\nabla \lambda\rangle dM \leq 0$. Then $X$ is parallel vector field if any one of the following conditions holds
\begin{enumerate}
\item $\displaystyle \int_M \mathrm{Ric}(X,X)\, dM 
\leq \frac{2}{m} \int_M |X|^2\mathrm{div}\,X \, dM$;

\item $X$ is conformal and $\displaystyle \int_M \mathrm{Ric}(X,X)dM\leq 0$;

\item $|X|$ is constant and $\displaystyle \int_M \mathrm{Ric}(X,X)dM\leq 0$.
\end{enumerate}
\end{corollary}

\begin{proof}
Integrating the expression of Lemma \ref{lem:laplacianX}, we obtain
\begin{equation*}
\int_M |\nabla X|^2 dM
= \int_M \mathrm{Ric}(X,X) dM
- \frac{2}{m}\int_M |X|^2 \,\mathrm{div}\,X \,\,dM
+(n-2)\int_M \langle X,\nabla \lambda\rangle dM.
\end{equation*}

In each of the cases above, the right-hand side is nonpositive. Hence,
$
\int_M |\nabla X|^2dM \leq 0,
$ 
which implies $|\nabla X|=0$, that is, $X$ is parallel.
\end{proof}

In the $m$-quasi-Einstein setting, Bahuaud et al.~\cite{Bahuaud2024} proved rigidity results under the assumption $|X|^2 = -m\lambda > 0$ and for $2 < m \leq 4$, showing in particular that $X$ is parallel. This was later improved by the first and third named authors~\cite{CarvalhoCosta}, who showed that the condition $|X|^2 = -m\lambda$ alone already implies that $X$ is parallel. We next study the corresponding statement in the generalized case.

\begin{proposition}\label{prop:parallel}
Let $(M^n,g,X,\lambda)$ be a closed generalized $m$-quasi-Einstein manifold such that
$
|X|^2 = -m\lambda.
$
Assume that one of the following conditions holds:
\begin{enumerate}
\item $m = 2(4-n)$;

\item $\displaystyle \frac{m}{2} + n - 4 > 0$ and 
$ 
\int_M \langle X, \nabla \lambda \rangle \, dM \leq 0;
$

\item $\displaystyle \frac{m}{2} + n - 4 < 0$ and 
$
\int_M \langle X, \nabla \lambda \rangle \, dM \geq 0.
$
\end{enumerate}
Then $X$ is parallel.
\end{proposition}
\begin{proof}
Proceeding as before, integrating the expression of Lemma \ref{lem:laplacianX}, we have that
\begin{equation*}
\int_M |\nabla X|^2 \, dM
= \int_M \mathrm{Ric}(X,X)\, dM
- \frac{2}{m}\int_M |X|^2 \,\mathrm{div}\,X \, dM
+ (n-2)\int_M \langle X,\nabla \lambda\rangle \, dM.
\end{equation*}

Combining \eqref{eq:quasi-einstein} with the identity $|X|^2 = -m\lambda$, we obtain
\begin{equation*}
\mathrm{Ric}(X,X) = \lambda |X|^2 + \frac{1}{m}|X|^4 
- \frac{1}{2}(\mathcal{L}_X g)(X,X)
= - \frac{1}{2}(\mathcal{L}_X g)(X,X).
\end{equation*}

Moreover,
$ 
(\mathcal{L}_X g)(X,X) = X(|X|^2) = -m\,\langle X,\nabla \lambda\rangle,
$ 
and hence
$ 
\mathrm{Ric}(X,X) = \frac{m}{2}\,\langle X,\nabla \lambda\rangle.
$ 

On the other hand, using the relation $|X|^2 = -m\lambda$  again  and the divergence theorem, we get
\begin{align*}
\int_M |X|^2 \,\mathrm{div}\,X \, dM
&= -m \int_M \lambda \,\mathrm{div}\,X \, dM \\
&= m \int_M \langle \nabla \lambda, X \rangle \, dM.
\end{align*}

Substituting these identities into the previous expression gives
\begin{equation*}
\int_M |\nabla X|^2 \, dM
= \frac{m}{2}\int_M \langle X,\nabla \lambda\rangle \, dM
- 2 \int_M \langle X,\nabla \lambda\rangle \, dM
+ (n-2)\int_M \langle X,\nabla \lambda\rangle \, dM.
\end{equation*}

Thus,
\begin{equation*}
\int_M |\nabla X|^2 \, dM
= \left(\frac{m}{2} + n - 4 \right)
\int_M \langle X,\nabla \lambda\rangle \, dM.
\end{equation*}

Hence, under all assumptions considered, we conclude that
$
\int_M |\nabla X|^2 \leq 0.
$
So, it follows that $|\nabla X|=0$, and therefore $X$ is parallel, which yields the result.
\end{proof}

We now recall some notions that will be useful in the sequel. According to \cite{StepanovShandra2003}, we say that a vector field $X$ on $(M^n,g)$ is an infinitesimal harmonic transformation if $\operatorname{tr}_g (\mathcal{L}_X\nabla)=0$, where $\nabla$ is the Levi-Civita connection of $M^n$. In turn, the definition of the rough Laplacian operator \(\bar{\Delta}\) acting on vector fields is given by \(\bar{\Delta}X=-\operatorname{tr}_g(\nabla^2X)\). Equivalently, if \(\{E_i\}\) is a local orthonormal frame field on $M^n$, then
\begin{equation}\label{eq:rough-Laplacian}
\bar{\Delta}X
=
\sum_i\left(
\nabla_{\nabla_{E_i}E_i}X
-\nabla_{E_i}\nabla_{E_i}X
\right).
\end{equation}

Next, we shall also denote by \(\operatorname{Ric}\) the \((1,1)\)-tensor associated with the Ricci tensor, that is,
$
\operatorname{Ric}(U,V)=\langle \operatorname{Ric}(U),V\rangle,
$
for all \(U,V\in \mathfrak{X}(M)\).

In \cite{StepanovShandra2003}, the authors proved that a vector field \(X\) generates an infinitesimal harmonic transformation on a Riemannian manifold \((M^{n},g)\) if and only if
$ 
\Delta X = 2\operatorname{Ric}(X),
$
where \(\Delta\) stands for the Laplacian related to the rough Laplacian through the Weitzenböck formula
$
\Delta X=\bar{\Delta}X+\operatorname{Ric}(X).
$

By direct computation, we can deduce that a Killing vector field is an example of an infinitesimal harmonic transformation.

We close this section by recalling the classical \emph{Yano formula}.
Let $(M^n,g)$ be a closed Riemannian manifold. 
Then, for any smooth vector field $Y$ on $M^n$, Yano’s identity (see \cite[p.~170]{Poor1981}) reads as
\begin{equation}\label{eq:yano}
\frac{1}{2}\int_M |\mathcal{L}_Y g|^{2}\,dM
=
\int_M \big(
|\nabla Y|^{2}
+(\operatorname{div}Y)^{2}
-\operatorname{Ric}(Y,Y)
\big)\,dM .
\end{equation}

\section{Proofs of The Results}\label{sec:proofs}

This section is dedicated to proving our main results.

\begin{proof}[Proof of Theorem A]

First, we may use Lemma~\ref{lem:laplacianX} to rewrite
the term $|\nabla X|^{2}-\operatorname{Ric}(X,X)$ appearing in \eqref{eq:yano}.
Indeed, we have
$$
|\nabla X|^{2}-\operatorname{Ric}(X,X)
=
\frac{1}{2}\,\Delta |X|^{2}
-\frac{2}{m}|X|^{2}\operatorname{div}X
-(2-n)\langle X,\nabla\lambda\rangle .
$$

Integrating \eqref{eq:yano} with \(Y=X\), and using the identity above, yields
$$
\frac{1}{2}\int_M |\mathcal{L}_X g|^{2}\,dM
=
\int_M \Big(
(\operatorname{div}X)^{2}
-\frac{2}{m}|X|^{2}\operatorname{div}X
-(2-n)\langle X,\nabla\lambda\rangle
+\frac{1}{2}\Delta |X|^{2}
\Big)\,dM .
$$

Because $M^n$ is closed, the integral of $\Delta |X|^{2}$ vanishes, and we finally obtain
\begin{equation}\label{eq27.2}
\frac{1}{2}\int_M |\mathcal{L}_X g|^{2}\,dM
=
\int_M \Big(
(\operatorname{div}X)^{2}
-\frac{2}{m}|X|^{2}\operatorname{div}X
-(2-n)\langle X,\nabla\lambda\rangle
\Big)\,dM .
\end{equation}

Applying the gradient to equation~\eqref{eq2}, one has
$$
\nabla S+\nabla(\operatorname{div}X)
=n\nabla\lambda +\frac{1}{m}\nabla(|X|^{2}).
$$

Taking the inner product with \(X\),
$$
\langle\nabla S, X\rangle
+\langle\nabla(\operatorname{div}X),X\rangle
=\frac{1}{m}\langle\nabla(|X|^{2}),X\rangle
+n\langle X,\nabla\lambda\rangle.
$$

Furthermore,
$$
\operatorname{div}(|X|^{2}X)
=|X|^{2}\operatorname{div}X
+2\langle\nabla_{X}X,X\rangle,
$$
$$
\operatorname{div}\big((\operatorname{div}X)X\big)
=(\operatorname{div}X)^{2}
+\langle\nabla(\operatorname{div}X),X\rangle.
$$

Integrating $\operatorname{div}\big((\operatorname{div}X)X\big)$ over a closed manifold
$M^n$ and applying the divergence theorem,
$$
0=\int_M \Big(
(\operatorname{div}X)^{2}
+\frac{1}{m}\langle\nabla(|X|^{2}),X\rangle
+n\langle X,\nabla\lambda\rangle-\langle\nabla S, X\rangle 
\Big)\,dM .
$$

Similarly, integrating \(\operatorname{div}(|X|^{2}X)\),
$$
\int_M \langle\nabla(|X|^{2}),X\rangle\,dM
=-\int_M |X|^{2}\operatorname{div}X\,dM .
$$

Consequently, we arrive at
\begin{equation}\label{eq27}
\int_M |X|^{2}\operatorname{div}X\,dM
=
m\int_M ((\operatorname{div}X)^{2}\
+n\langle X,\nabla\lambda\rangle-\langle\nabla S, X\rangle)\,dM .
\end{equation}


Combining \eqref{eq27} and \eqref{eq27.2},
we obtain
$$
\begin{aligned}
\frac{1}{2}\int_M |\mathcal{L}_X g|^{2}\,dM
&=
\int_M \Big(
(\operatorname{div}X)^{2}
-2(\operatorname{div}X)^{2}
-2n\langle X,\nabla\lambda\rangle
-(2-n)\langle X,\nabla\lambda\rangle
+2\langle\nabla S, X\rangle)\Big)\,dM \\
&=
-\int_M \Big(
(\operatorname{div}X)^{2}
+(n+2)\langle X,\nabla\lambda\rangle
-2\langle\nabla S, X\rangle)\Big)\,dM .
\end{aligned}
$$

The previous identity and the assumptions $\int_M \langle X,\nabla\lambda\rangle\ge0$, and $\int_M \langle X,\nabla S\rangle\le0$ gives
$$
\int_M |\mathcal{L}_X g|^{2}\,dM \le 0.
$$

Hence, $\mathcal{L}_X g=0$, and $X$ is Killing. So, the proof is completed.
\end{proof}

We now turn to the proof of Proposition \ref{prop-div}.
\begin{proof}[Proof of Proposition B]Assuming that $\operatorname{div} X = 0$,
we apply Lemma~\ref{lem:laplacianX} to obtain
$$
\frac{1}{2}\,\Delta |X|^{2}
= |\nabla X|^{2}
- \operatorname{Ric}(X,X)
+ (2-n)\,\langle X,\nabla \lambda\rangle .
$$

Integrating this identity over $M^n$ and comparing it with the formula
\eqref{eq:yano}, we deduce that
$$
\frac{1}{2}\int_M | \mathcal{L}_X g |^2 \, dM = (n-2) \int_M\,\langle X,\nabla \lambda\rangle dM.
$$

Since \(\operatorname{div}X=0\), the divergence theorem implies
$$
0 = \int_M \operatorname{div}(\lambda X)\, dM
= \int_M \langle X,\nabla \lambda \rangle \, dM .
$$
Consequently, \(\mathcal{L}_X g =0\), so \(X\) is a Killing vector field.
\end{proof}

We are now in a position to prove Corollary~E.

\begin{proof}[Proof of Corollary E]

Evaluating \eqref{eq:quasi-einstein} along the pair \((X,X)\),
$$
\operatorname{Ric}(X,X) + \langle \nabla_X X, X \rangle
= \frac{1}{m}\lvert X\rvert^4 + \lambda \lvert X\rvert^2 .
$$
On the other hand,
$$
\operatorname{div}(|X|^2 X)
= \langle \nabla |X|^2, X \rangle + |X|^2 \operatorname{div} X,
$$
Since \(\operatorname{div}X=0\), integration over \(M^n\) together with the divergence theorem gives
$$
\int_M \langle \nabla_X X, X \rangle \, dM = 0 .
$$
Therefore,
$$
\int_M \operatorname{Ric}(X,X) \, dM
= \int_M\left(\frac{1}{m}\lvert X\rvert^4 + \lambda \lvert X\rvert^2\right) dM .
$$

Under assumption \(\operatorname{div}X=0\), Lemma~\ref{lem:laplacianX} reduces to
$$
\frac{1}{2}\,\Delta |X|^{2}
= |\nabla X|^{2}
- \operatorname{Ric}(X,X)
+ (2-n)\,\langle X,\nabla \lambda\rangle .
$$

Integrating over $M^n$ and using the divergence theorem
$$
\int_M \left(
|\nabla X|^2
- \operatorname{Ric}(X,X)
+ (2-n)\,\langle X,\nabla \lambda\rangle
\right) dM = 0 .
$$

Combining this identity with the previous integral formula for
$\int_M \operatorname{Ric}(X,X)\, dM$, we conclude that
$$
\int_M \left(
|\nabla X|^2
- \frac{1}{m}\lvert X\rvert^4
- \lambda \lvert X\rvert^2
+ (2-n)\,\langle X,\nabla \lambda\rangle
\right) dM = 0 .
$$

Using again \(\operatorname{div}X=0\), integration by parts gives
$$
0 = \int_M \operatorname{div}(\lambda X)\, dM
= \int_M \langle X,\nabla \lambda \rangle \, dM .
$$

Hence,
$$
\int_M
|\nabla X|^2 dM=\int_M
\frac{1}{m}\lvert X\rvert^4
 +\lambda \lvert X\rvert^2
 dM  .
$$

Under assumptions $m<0$ and $\lambda \le 0$, the right-hand side is nonpositive,
and therefore $\int_M
|\nabla X|^2dM$  must vanish identically. This implies $X =0$ on $M^n$.
\end{proof}

We now prove the following theorem.

\begin{proof}[Proof of Theorem F]
The generalized $m$-quasi-Einstein equation \eqref{eq:quasi-einstein} can be written,
for any vector fields $U,V$ on $M^n$, as
\begin{equation}\label{eq:gen-mqe}
\langle \nabla_U X, V \rangle + \langle \nabla_V X, U \rangle
+ 2\,\operatorname{Ric}(U,V)
- \frac{2}{m}\, \langle U,X\rangle \langle V,X\rangle
= 2\lambda\, \langle U,V\rangle .
\end{equation}

Setting $U=X$ in \eqref{eq:gen-mqe} and using the hypothesis
$\nabla_{X}X=0$, we obtain 

\begin{equation}\label{eq:step1}
\langle \nabla_V X, X \rangle
+ 2\,\langle \operatorname{Ric}(X), V \rangle
- \frac{2}{m}\lvert X\rvert^{2} \langle V,X\rangle
= 2\lambda\, \langle X,V\rangle.
\end{equation}

Since
$$
\langle \nabla_V X, X \rangle = \frac{1}{2}\, V(\lvert X\rvert^{2}),
$$
and the Riemannian metric is non-degenerate, equation \eqref{eq:step1}
can be rewritten in vectorial form as
\begin{equation}\label{eq:step2}
\frac{1}{2}\nabla \lvert X\rvert^{2} + 2\operatorname{Ric}(X)
= 2\left(\lambda + \frac{\lvert X\rvert^{2}}{m}\right) X.
\end{equation}

Let $\{e_i\}_{i=1}^n$ be a local orthonormal frame on $M^n$.
Differentiating \eqref{eq:step2} with respect to $e_i$,
taking the inner product with $e_i$, and summing over $i$,
we obtain
\begin{equation}\label{eq:step3}
\frac{1}{2}\Delta |X|^{2}
+ 2\,\operatorname{div}(\operatorname{Ric}(X))
= 2\left(\lambda + \frac{|X|^{2}}{m}\right)\operatorname{div} X
+ 2\langle\nabla\lambda, X\rangle + \frac{4}{m}\,\langle \nabla_{X}X, X\rangle .
\end{equation}

By assumption $\nabla_{X}X=0$, the last term vanishes.
Integrating \eqref{eq:step3} over the closed manifold $M^n$ and applying
the divergence theorem, together with the identity
$$
2\lambda\,\operatorname{div} X
+ 2\langle X, \nabla\lambda\rangle
= 2\,\operatorname{div}(\lambda X),
$$
we obtain

$$
\int_{M} |X|^{2}\, \operatorname{div} X \, dM = 0.
$$

Therefore, we use Lemma~\ref{lem:laplacianX} to infer
$$
\int_M \left(
|\nabla X|^{2}
- \operatorname{Ric}(X,X)
+ (2-n)\,\langle X,\nabla \lambda\rangle
\right) dM = 0 .
$$

Since $\int_M \operatorname{Ric}(X,X) dM \le 0$ and
$\int_M \langle X,\nabla \lambda\rangle\, dM \le 0$, we conclude that
$\nabla X =0$. Thus, $\operatorname{div} X =0$ and $|X|$ is constant. Thus, by Lemma~\ref{lem:laplacianX}, we have
$\operatorname{Ric}(X,X) = (2-n)\langle X,\nabla\lambda\rangle$.
Together with
$\int_M \operatorname{Ric}(X,X)\,dM \le 0$ and
$\int_M \langle X,\nabla\lambda\rangle\,dM \le 0$,
this implies $\operatorname{Ric}(X,X)=0$.
Equation~\eqref{eq:quasi-einstein} then gives
$
0=\frac{1}{m}|X|^{4}+\lambda|X|^{2}.
$

Therefore, if $m$ and $\lambda$ have the same sign, then $X=0$.
If $\lambda=0$, then $X=0$ and $(M^n,g)$ is Ricci-flat.
If \(X\not=0\), then from
$
0=\frac{1}{m}|X|^{4}+\lambda|X|^{2},
$
it follows that
$
|X|^{2}=-m\lambda
$
on the set \(\{X\neq0\}\). This finishes the proof of Theorem. 
\end{proof}

We now present the identity below, which is interesting in their own right and will play an important role in the proof of Theorem \ref{teo6.1}.

\begin{lemma}\label{eq10}
Let $(M^n,g,X,\lambda)$ be a generalized quasi-Einstein manifold. Then, for every smooth vector field $Z$ on $M$, the following identity hold:
\begin{align}
\langle 2\Ric X - \Delta X,Z\rangle
&=
\frac{2}{m}\Bigl(
\operatorname{div}(X)\,\langle Z,X\rangle
+\langle \nabla_XX,Z\rangle
-\langle \nabla_ZX,X\rangle
\Bigr) \notag\\
&\qquad
-(n-2)\langle \nabla\lambda,Z\rangle, \label{neweq10.1}
\end{align}
\end{lemma}


\begin{proof} The identity \eqref{neweq10.1} can be viewed as an extension of Equation~(5.4) in Ghosh~\cite{ghosh} to the generalized setting. Although the proof follows essentially the same lines and also relies on Yano's commutation formula, we include it here for completeness and convenience of the reader. We also point out that there is a minor typographical error in the corresponding Yano's commutation formula in~\cite{ghosh}.

We first derive an expression for $\mathcal{L}_X\nabla$ that will be used later. Differentiating \eqref{eq:quasi-einstein} along an arbitrary vector field $Z$, we obtain
\begin{align}
(\nabla_Z\mathcal{L}_Xg)(U,V)
+2(\nabla_Z\operatorname{Ric})(U,V)
&=
\frac{2}{m}\Bigl(
\langle U,\nabla_ZX\rangle \langle V,X\rangle
\notag\\
&\qquad\qquad
+\langle U,X\rangle \langle V,\nabla_ZX\rangle
\Bigr)
\notag\\
&\quad
+2(Z\lambda)\langle U,V\rangle .
\end{align}

Using the commutation formula (see Yano \cite{Yan70})
\[
(\mathcal{L}_X\nabla_Z g-\nabla_Z\mathcal{L}_X g-\nabla_{[X,Z]}g)(U,V)
=
-\inner{(\mathcal{L}_X\nabla)(Z,U)}{V}
-\inner{(\mathcal{L}_X\nabla)(Z,V)}{U},
\]
we deduce
\begin{equation}\label{equation3}
\begin{aligned}
&
\langle (\mathcal{L}_X\nabla)(Z,U),V\rangle
+
\langle (\mathcal{L}_X\nabla)(Z,V),U\rangle
\\
&\qquad
=
-2(\nabla_Z\operatorname{Ric})(U,V)
+2(Z\lambda)\langle U,V\rangle
\\
&\qquad\quad
+\frac{2}{m}\Bigl[
\langle U,\nabla_ZX\rangle \langle V,X\rangle
+
\langle U,X\rangle \langle V,\nabla_ZX\rangle
\Bigr].
\end{aligned}
\end{equation}

Do a cyclic permutation of \(U,V,Z\) in equation~\eqref{equation3} to obtain the following two additional identities:
\begin{equation}\label{equation4}
\begin{split}
&
\left\langle (\mathcal{L}_X\nabla)(U,V),Z \right\rangle
+
\left\langle (\mathcal{L}_X\nabla)(U,Z),V \right\rangle
\\
&\qquad
=
-2(\nabla_U\operatorname{Ric})(V,Z)
+2(U\lambda)\left\langle V,Z \right\rangle
\\
&\qquad\quad
+\frac{2}{m}\Bigl[
\left\langle V,\nabla_UX \right\rangle
\left\langle Z,X \right\rangle
+
\left\langle V,X \right\rangle
\left\langle Z,\nabla_UX \right\rangle
\Bigr].
\end{split}
\end{equation}

Similarly, we obtain
\begin{equation}\label{equation5}
\begin{split}
&
\left\langle (\mathcal{L}_X\nabla)(V,Z),U \right\rangle
+
\left\langle (\mathcal{L}_X\nabla)(V,U),Z \right\rangle
\\
&\qquad
=
-2(\nabla_V\operatorname{Ric})(Z,U)
+2(V\lambda)\left\langle Z,U \right\rangle
\\
&\qquad\quad
+\frac{2}{m}\Bigl[
\left\langle Z,\nabla_VX \right\rangle
\left\langle U,X \right\rangle
+
\left\langle Z,X \right\rangle
\left\langle U,\nabla_VX \right\rangle
\Bigr].
\end{split}
\end{equation}

Next,  taking into account that \(\mathcal{L}_X\nabla\) is symmetric, adding \eqref{equation4} and \eqref{equation5}, and then subtracting \eqref{equation3}, we obtain the following
\begin{equation}\label{eq123}
\begin{split}
\left\langle (\mathcal{L}_X\nabla)(U,V),Z \right\rangle
&=(\nabla_Z\operatorname{Ric})(U,V)
-(\nabla_V\operatorname{Ric})(Z,U)
-(\nabla_U\operatorname{Ric})(V,Z)
\\
&\quad
+\frac1m\Big[
\left\langle Z,X \right\rangle
\Big(
\left\langle \nabla_UX,V \right\rangle
+\left\langle \nabla_VX,U \right\rangle
\Big)
\\
&\qquad\quad
+\left\langle V,X \right\rangle
\Big(
\left\langle \nabla_UX,Z \right\rangle
-\left\langle \nabla_ZX,U \right\rangle
\Big)
\\
&\qquad\quad
+\left\langle U,X \right\rangle
\Big(
\left\langle \nabla_VX,Z \right\rangle
-\left\langle \nabla_ZX,V \right\rangle
\Big)
\Big]
\\
&\quad
-(Z\lambda)\left\langle U,V \right\rangle
+(V\lambda)\left\langle Z,U \right\rangle
+(U\lambda)\left\langle V,Z \right\rangle .
\end{split}
\end{equation}

Thus, taking \(U=V=E_i\), where \(\{E_1,\dots,E_n\}\) is a local orthonormal frame field, and summing over \(i\), the right-hand side of \eqref{eq123} becomes
\[
\frac{2}{m}\left\{
\operatorname{div}(X)\langle Z,X\rangle
+\langle \nabla_XX,Z\rangle
-\langle \nabla_ZX,X\rangle
\right\}
-(n-2)\langle\nabla\lambda,Z\rangle .
\]

Now, let $R$ denote the Riemann curvature tensor of $g$. 
By the well-known identity due to Yano~\cite{Yan70} (see Equation~(5.13), p.~23),

\[
\left\langle (\mathcal{L}_X\nabla)(U,V),Z \right\rangle
=
\left\langle
\nabla_U\nabla_VX
-\nabla_{\nabla_UV}X
+R(X,U)V,
\, Z
\right\rangle,
\]


Hence,% taking \(U=V=E_i\), where \(\{E_1,\dots,E_n\}\) is a local orthonormal frame field, and summing over \(i\), we obtain
\[
\sum_{i=1}^n
\left\langle
(\mathcal{L}_X\nabla)(E_i,E_i),Z
\right\rangle
=
\sum_{i=1}^n
\left\langle
\nabla_{E_i}\nabla_{E_i}X
-\nabla_{\nabla_{E_i}E_i}X
+R(X,E_i){E_i},
Z
\right\rangle .
\]

It follows immediately from the definitions \eqref{eq:rough-Laplacian} that the left-hand side of \eqref{eq123} may be rewritten as
$
\left\langle
-\bar{\Delta}X+\Ric X,Z
\right\rangle .
$

Using the Weitzenb\"ock formula
\(
\Delta X=\bar{\Delta}X+\Ric X,
\)
it follows that
$
\langle 2\Ric X - \Delta X,Z\rangle
=
\frac{2}{m}\left\{
\operatorname{div}(X)\langle Z,X\rangle
+\langle \nabla_XX,Z\rangle
-\langle \nabla_ZX,X\rangle
\right\}
-(n-2)\langle\nabla\lambda,Z\rangle,
$ as we wished to prove.


\end{proof}

Now, we are in position to prove Theorem \ref{teo6.1}.

\begin{proof}[Proof of Theorem C]

Since \(X\) is divergence-free, Proposition~\ref{prop-div} shows that \(X\) is a Killing vector field. Therefore, equation~\eqref{neweq10.1} in Lemma~\ref{eq10} yields
$ 
0
=
\frac{4}{m}\langle \nabla_XX,Z\rangle
-(n-2)\langle \nabla\lambda,Z\rangle,
$
which implies
\begin{equation}\label{eq104}
\nabla_XX
=
\frac{m}{4}(n-2)\nabla\lambda.
\end{equation}

Moreover, taking the divergence in \eqref{eq:quasi-einstein} and using the contracted second Bianchi identity
$\nabla S = 2\operatorname{div}\operatorname{Ric},
$
we obtain
\begin{equation}\label{eq:div}
\nabla S
=
2\nabla\lambda
+
\frac{2}{m}\nabla_XX,
\end{equation}
where we used that \(X\) is Killing. Combining \eqref{eq104} and \eqref{eq:div}, we obtain
$$
\nabla S
=
2\nabla\lambda
+
\frac{n-2}{2}\nabla\lambda
=
\frac{n+2}{2}\nabla\lambda.
$$

This completes the proof.
\end{proof}

We close our paper by presenting the proof of Theorem \ref{D}.

\begin{proof}[Proof of Theorem D]
Using that $X$ is an infinitesimal harmonic transformation and assuming either $n=2$ or that $\lambda$ is invariant along the flow of $X$, equation \eqref{neweq10.1} reduces to
\begin{equation*}
\frac{2}{m}\Bigl(
\operatorname{div}(X)\,\langle Z,X\rangle
+\langle \nabla_XX,Z\rangle
-\langle \nabla_ZX,X\rangle
\Bigr)=0.
\end{equation*}

Setting $Z=X$ in the above identity gives
$\operatorname{div}(X)\,|X|^2=0.
$ 
The rest of the argument proceeds exactly as in the proof of Theorem 5.1 of Ghosh \cite{ghosh}.

Now, assume that $X$ is Killing. Since $X$ is geodesic, it follows that $|X|$ is constant. On the other hand, using again that $X$ is a Killing vector field, the equation \eqref{neweq10.1} becomes
$\frac{4}{m}\langle \nabla_XX,Z\rangle
=(n-2)\langle \nabla\lambda,Z\rangle.
$
Since this identity holds for every vector field $Z$, we obtain
$ 
\frac{4}{m}\nabla_XX=(n-2)\nabla\lambda.
$ 
As $X$ is geodesic, we infer
$ 
(n-2)\nabla\lambda=0.
$
Then, if $n\geq 3$, we have that $\lambda$ is constant, which proves the result.
\end{proof}

After finishing this work, we became aware of the article \cite{GulerDe2022}.
In that paper, the authors state the following result:

\begin{lemma}\label{lem:gulerde}
Let $(M^n,g,X,\lambda)$ be a generalized quasi-Einstein manifold. Then
\begin{equation}\label{eq:gulerde-lemma1}
- \langle\bar{\Delta} X, Z\rangle + \operatorname{Ric}(X, Z)
= 2\Bigl(\, \langle X, Z\rangle\, \operatorname{div} X
+ \langle\nabla_X X, Z\rangle
- \langle\nabla_Z X, X \rangle\Bigr)
+ (n - 2)\, d\lambda(Z),
\end{equation}
for all $Z \in \mathfrak{X}(M)$.
\end{lemma}
We note that identity \eqref{eq:gulerde-lemma1} differs from our formula
\eqref{neweq10.1} by the sing on the term $\, \langle\nabla \lambda, Z\rangle$.
The discrepancy appears to arise from a sign error
in Equation~(24) of \cite{GulerDe2022}, which propagates into the subsequent computations. We also note that, in their setting, the parameter \(m\) is assumed to be \(1\).

Moreover, the authors use Lemma~\ref{lem:gulerde} to derive the following result:

\begin{theorem}[{\cite[Theorem 2]{GulerDe2022}}]\label{thm:gulerde}
Let $(M^n,g,X,\lambda)$ be a generalized quasi-Einstein manifold such that
the associated $1$-form of the potential vector field $X$ is harmonic.
Then $\Delta X = 4 QX$.
\end{theorem}

Since the proof of Theorem~\ref{thm:gulerde} depends directly on
Lemma~\ref{lem:gulerde}, the correction of the latter also affects the statement
of this theorem. More precisely, the coefficient in the identity for $\Delta X$
must be replaced by $0$ instead of $4,$ that is, $\Delta X =0$.

We also point out a further issue in \cite{GulerDe2022}. More precisely, the statement of Theorem 3 does  not hold in full generality. The authors state:
\begin{theorem}[{\cite[Theorem 3]{GulerDe2022}}]
Let $(M^n,g,V,\lambda)$ be a closed generalized quasi-Einstein manifold whose potential vector field $V$ is parallel. Then $V=0$.
\end{theorem}

Indeed, if $V$ is parallel, i.e.,
$
\nabla V=0,
$
then $|V|$ is constant and $\operatorname{div}V=0$. Hence, Lemma~\ref{lem:laplacianX} immediately yields equation (44) of \cite{GulerDe2022},
$
\operatorname{Ric}(V,V)+(n-2)\langle\nabla\lambda,V\rangle=0,
$
without the need to integrate or use the compactness assumption.
On the other hand, equation (45) in \cite{GulerDe2022} asserts that
$$
\operatorname{Ric}(X,Y)
=
V^\flat(X)V^\flat(Y)
-
|V|^2 \langle X,Y \rangle,
$$
for all $X,Y\in\mathfrak{X}(M)$. However, this identity does \emph{not} imply that $M^n$ is Ricci-flat. It only shows that $V$ belongs to the nullity of the Ricci tensor, since we have in this case,
$ 
\operatorname{Ric}(V,X)=0
$
for all $X\in\mathfrak{X}(M)$.

A similar issue also appears in \cite[Equation (2.6), Lemma 1 item (3)]{BarrosRibeiro2012}, and was already pointed out by the first and third authors in \cite{CarvalhoCosta}. A counterexample is given by the Riemannian product
$
(\mathbb{S}^1\times N^{n-1},\, d\theta^2+g_N),
$
where $(N^{n-1},g_N)$ is an Einstein manifold.
In this case, the vector field
$
X=\pm\sqrt{-m\lambda}\,\partial_\theta, \,\,\,  m\lambda<0,
$
is parallel and satisfies
$
\operatorname{Ric}_X^m=\lambda g,
$
while clearly
$ 
X\neq0.
$  Moreover, we show in \cite{CarvalhoCosta} that the above counterexample is the only nontrivial one.
Therefore, the conclusion of Theorem 3 does not hold in full generality.

\textbf{Declarations}\\

\textbf{Data Availability} This manuscript has no associated data.\\

\textbf{Conflict of interest} The authors declare that there is no  conflict of interest.\\

\textbf{Acknowledgments} The authors would like to thank A. Ghosh for his kind interest in this work, for clarifying several aspects of his article, which greatly improved the manuscript. We are also grateful to I.~Domingos for his kind interest in this work.
The first author thanks the Federal University of Alagoas (UFAL), Arapiraca campus, where part of this work was developed during a research visit. The second author thanks the Department of Mathematics at the Federal University of Pernambuco for its hospitality during his visit, where part of this work was carried out. The third author was supported by CNPq/Brazil Grant 409513/2023-7.
\bibliographystyle{plain}
\bibliography{references}
\end{document}

